\documentclass{article}
\usepackage{amsmath,amssymb,amsthm}
\usepackage{hyperref}

\newtheorem{theorem}{Theorem}
\newtheorem{lemma}[theorem]{Lemma}
\newtheorem{proposition}[theorem]{Proposition}
\newtheorem{corollary}[theorem]{Corollary}
\theoremstyle{definition}
\newtheorem{definition}[theorem]{Definition}
\newtheorem{remark}[theorem]{Remark}

\title{On a cubic identity in hyperfields}
\author{}
\date{}

\begin{document}

\maketitle

\begin{abstract}
We prove that in any hyperfield, the ternary additive relation
\(0\in x\boxplus y\boxplus z\) implies the cubic relation
\[
0\in x^3\boxplus y^3\boxplus z^3\boxplus(-xyz)\boxplus(-xyz)\boxplus(-xyz).
\]
The proof uses only the hyperfield axioms and a standard factorization identity.
\end{abstract}

\section{Hyperfields and conventions}

\begin{definition}[Hyperfield]
A \emph{hyperfield} is a set \(\mathbb H\) equipped with:
\begin{enumerate}
    \item a commutative, associative multivalued addition
    \[
    \boxplus:\mathbb H\times \mathbb H\to \mathcal P(\mathbb H)\setminus\{\varnothing\},
    \]
    with neutral element \(0\), such that for every \(a\in\mathbb H\) there is a unique
    \(-a\in\mathbb H\) with \(0\in a\boxplus(-a)\), and such that
    \(a\in b\boxplus c\) implies \(c\in(-b)\boxplus a\);
    \item a multiplication \(\cdot\) making \(\mathbb H^\times=\mathbb H\setminus\{0\}\)
    an abelian group, with \(0\cdot a=a\cdot 0=0\);
    \item distributivity:
    \[
    a\cdot(b\boxplus c)=ab\boxplus ac
    \]
    for all \(a,b,c\in\mathbb H\).
\end{enumerate}
We write sums without parentheses, using associativity of \(\boxplus\). Thus
\(x\boxplus y\boxplus z\) denotes the set obtained by adding \(x,y,z\) in any order.
\end{definition}

Throughout, all products are written multiplicatively, and we use the standard
identities
\[
(x+y)^2=x^2\boxplus 2xy\boxplus y^2,\qquad
(x+y)^3=x^3\boxplus 3x^2y\boxplus 3xy^2\boxplus y^3,
\]
which follow from distributivity and associativity. Here \(2xy\) means
\(xy\boxplus xy\), and \(3t\) means \(t\boxplus t\boxplus t\).

\section{The main result}

\begin{lemma}\label{lem:factorization}
For any elements \(a,b,c\) of a hyperfield, the following two multivalued sums are equal
as subsets of \(\mathbb H\):
\[
a^3\boxplus b^3\boxplus c^3\boxplus(-abc)\boxplus(-abc)\boxplus(-abc)
=
(a+b+c)\Bigl(a^2\boxplus b^2\boxplus c^2\boxplus(-ab)\boxplus(-bc)\boxplus(-ca)\Bigr).
\]
\end{lemma}

\begin{proof}
Expand the right-hand side using distributivity. The product of
\((a+b+c)\) with
\(a^2+b^2+c^2-ab-bc-ca\) gives the eighteen terms
\[
\begin{aligned}
& a^3+ab^2+ac^2-a^2b-abc-a^2c\\
&\quad +a^2b+b^3+bc^2-ab^2-b^2c-abc\\
&\quad +a^2c+b^2c+c^3-abc-bc^2-ac^2 .
\end{aligned}
\]
The mixed terms appear in pairs with opposite signs. For example, \(ab^2\) and
\(-ab^2\) occur together, and since \(ab^2\in ab^2\boxplus(-ab^2)\) is not the issue:
in the hyperfield sum, the set \(ab^2\boxplus(-ab^2)\) contains \(0\). Grouping each
such pair and using associativity of \(\boxplus\), the total sum contains exactly
\[
a^3\boxplus b^3\boxplus c^3\boxplus(-abc)\boxplus(-abc)\boxplus(-abc).
\]
Conversely, every term of the right-hand side is a sum of terms appearing in the
left-hand side. Hence the two sets are equal. This is the usual symmetric
factorization
\[
a^3+b^3+c^3-3abc=(a+b+c)(a^2+b^2+c^2-ab-bc-ca),
\]
read in the hyperfield setting.
\end{proof}

\begin{theorem}\label{thm:main}
Let \(\mathbb H\) be a hyperfield and let \(x,y,z\in\mathbb H\). If
\[
0\in x\boxplus y\boxplus z,
\]
then
\[
0\in x^3\boxplus y^3\boxplus z^3\boxplus(-xyz)\boxplus(-xyz)\boxplus(-xyz).
\]
\end{theorem}

\begin{proof}
Set
\[
S=x\boxplus y\boxplus z.
\]
By hypothesis, \(0\in S\). Let
\[
Q=xy\boxplus yz\boxplus zx.
\]
Consider the element
\[
F=x^2\boxplus y^2\boxplus z^2\boxplus(-xy)\boxplus(-yz)\boxplus(-zx).
\]
By Lemma~\ref{lem:factorization},
\[
x^3\boxplus y^3\boxplus z^3\boxplus(-xyz)\boxplus(-xyz)\boxplus(-xyz)
=
S\cdot F.
\]
Since \(0\in S\), the product \(S\cdot F\) contains \(0\cdot f=0\) for every
\(f\in F\). In particular,
\[
0\in S\cdot F
=
x^3\boxplus y^3\boxplus z^3\boxplus(-xyz)\boxplus(-xyz)\boxplus(-xyz).
\]
This proves the desired implication.
\end{proof}

\begin{corollary}
In every hyperfield, the ternary relation
\[
0\in x\boxplus y\boxplus z
\]
implies the cubic relation
\[
0\in x^3\boxplus y^3\boxplus z^3\boxplus(-xyz)\boxplus(-xyz)\boxplus(-xyz).
\]
\end{corollary}

\begin{remark}
The theorem is a hyperfield analogue of the classical identity
\[
x+y+z=0\quad\Longrightarrow\quad x^3+y^3+z^3-3xyz=0.
\]
The proof above uses only the hyperfield axioms: associativity and commutativity of
\(\boxplus\), existence and uniqueness of additive inverses, and distributivity of
multiplication over multivalued addition.
\end{remark}

\end{document}